\section{Problem and data description}
\label{sec:data}

% Dataset numbers below come from results/a2/eda/suim/summary.json through
% scripts/report/a2_data_tables.py; regenerate instead of editing them.
% Generated by scripts/report/a2_data_tables.py from results/a2/eda/suim/summary.json. Do not edit.
\newcommand{\suimTrainValOnDisk}{1{,}525}
\newcommand{\suimTestOnDisk}{110}
\newcommand{\suimOnDiskTotal}{1{,}635}
\newcommand{\suimUnpaired}{0}
\newcommand{\suimDecodeErrors}{0}
\newcommand{\suimSizeMismatches}{37}
\newcommand{\suimOffCodeMasks}{85}
\newcommand{\suimTestOffCodeMasks}{0}
\newcommand{\suimOffCodeShareMin}{2.6}
\newcommand{\suimOffCodeShareMax}{85}
\newcommand{\suimQualityPassed}{1{,}440}
\newcommand{\suimTestDuplicatesDropped}{10}
\newcommand{\suimTrainValUsable}{1{,}430}
\newcommand{\suimTestUsable}{110}
\newcommand{\suimUsableTotal}{1{,}540}
\newcommand{\suimStorage}{1.8}
\newcommand{\suimImageStorage}{0.2}
\newcommand{\suimMaskStorage}{1.7}
\newcommand{\suimForegroundClasses}{7}
\newcommand{\suimResolutionCount}{11}
\newcommand{\suimDominantResolution}{640$\times$480}
\newcommand{\suimDominantResolutionShare}{84}
\newcommand{\suimHashBits}{64}
\newcommand{\suimHashDistance}{12}
\newcommand{\suimPixelTolerance}{5}
\newcommand{\suimHashCandidates}{2{,}035}
\newcommand{\suimConfirmedPairs}{76}
\newcommand{\suimDuplicateGroups}{65}
\newcommand{\suimCrossSplitGroups}{10}
\newcommand{\suimWithinGroups}{55}
\newcommand{\suimWithinGroupImages}{115}
\newcommand{\suimSplitGroups}{1{,}370}
\newcommand{\suimSplitSeed}{42}
\newcommand{\suimValFolds}{5}
\newcommand{\suimTrainImages}{1{,}143}
\newcommand{\suimValImages}{287}
\newcommand{\suimTrainPercent}{79.9}
\newcommand{\suimValPercent}{20.1}
\newcommand{\suimMaxShareGap}{2.9}
\newcommand{\suimShareTrainValBW}{30.99}
\newcommand{\suimShareTestBW}{41.81}
\newcommand{\suimMedianAreaBW}{33.6}
\newcommand{\suimImagesTrainValBW}{1{,}185}
\newcommand{\suimImagesTrainBW}{949}
\newcommand{\suimShareTrainBW}{31.4}
\newcommand{\suimImagesValBW}{236}
\newcommand{\suimShareValBW}{29.4}
\newcommand{\suimShareTrainValHD}{2.34}
\newcommand{\suimShareTestHD}{3.69}
\newcommand{\suimMedianAreaHD}{4.9}
\newcommand{\suimImagesTrainValHD}{359}
\newcommand{\suimImagesTrainHD}{287}
\newcommand{\suimShareTrainHD}{2.4}
\newcommand{\suimImagesValHD}{72}
\newcommand{\suimShareValHD}{2.1}
\newcommand{\suimShareTrainValPF}{2.29}
\newcommand{\suimShareTestPF}{3.08}
\newcommand{\suimMedianAreaPF}{7.4}
\newcommand{\suimImagesTrainValPF}{218}
\newcommand{\suimImagesTrainPF}{175}
\newcommand{\suimShareTrainPF}{2.2}
\newcommand{\suimImagesValPF}{43}
\newcommand{\suimShareValPF}{2.4}
\newcommand{\suimShareTrainValWR}{6.77}
\newcommand{\suimShareTestWR}{7.28}
\newcommand{\suimMedianAreaWR}{43.2}
\newcommand{\suimImagesTrainValWR}{245}
\newcommand{\suimImagesTrainWR}{197}
\newcommand{\suimShareTrainWR}{6.7}
\newcommand{\suimImagesValWR}{48}
\newcommand{\suimShareValWR}{7.0}
\newcommand{\suimShareTrainValRO}{0.56}
\newcommand{\suimShareTestRO}{0.90}
\newcommand{\suimMedianAreaRO}{3.8}
\newcommand{\suimImagesTrainValRO}{86}
\newcommand{\suimImagesTrainRO}{69}
\newcommand{\suimShareTrainRO}{0.5}
\newcommand{\suimImagesValRO}{17}
\newcommand{\suimShareValRO}{0.7}
\newcommand{\suimShareTrainValRI}{35.64}
\newcommand{\suimShareTestRI}{19.47}
\newcommand{\suimMedianAreaRI}{56.1}
\newcommand{\suimImagesTrainValRI}{964}
\newcommand{\suimImagesTrainRI}{761}
\newcommand{\suimShareTrainRI}{35.0}
\newcommand{\suimImagesValRI}{203}
\newcommand{\suimShareValRI}{37.9}
\newcommand{\suimShareTrainValFV}{7.26}
\newcommand{\suimShareTestFV}{6.77}
\newcommand{\suimMedianAreaFV}{7.2}
\newcommand{\suimImagesTrainValFV}{958}
\newcommand{\suimImagesTrainFV}{760}
\newcommand{\suimShareTrainFV}{7.3}
\newcommand{\suimImagesValFV}{198}
\newcommand{\suimShareValFV}{7.2}
\newcommand{\suimShareTrainValSR}{14.15}
\newcommand{\suimShareTestSR}{17.00}
\newcommand{\suimMedianAreaSR}{29.5}
\newcommand{\suimImagesTrainValSR}{571}
\newcommand{\suimImagesTrainSR}{458}
\newcommand{\suimShareTrainSR}{14.4}
\newcommand{\suimImagesValSR}{113}
\newcommand{\suimShareValSR}{13.2}
\newcommand{\suimOffCodeShareSameSize}{35}
\newcommand{\suimOffCodeShareTaller}{23}
\newcommand{\suimTallerExtraRows}{55}
\newcommand{\suimDuplicateScoreLowest}{0.0}
\newcommand{\suimDuplicateScoreHighest}{0.5}
\newcommand{\suimImbalanceRatio}{64}
\newcommand{\suimInputHeight}{240}
\newcommand{\suimInputWidth}{320}
\newcommand{\suimMaskThreshold}{128}
\newcommand{\suimNormMean}{(0.485, 0.456, 0.406)}
\newcommand{\suimNormStd}{(0.229, 0.224, 0.225)}
\newcommand{\suimExampleBatchSize}{4}
\newcommand{\suimExampleFlipped}{1}
\newcommand{\suimExampleClasses}{8}
\newcommand{\suimExampleValueMin}{-2.12}
\newcommand{\suimExampleValueMax}{2.64}
\newcommand{\suimBatch}{4}
\newcommand{\suimAccumulation}{4}
\newcommand{\suimEffectiveBatch}{16}


\subsection{Task and research question}
\label{sec:task}
Autonomous underwater vehicles and remotely operated robots rely on cameras to
follow divers, avoid obstacles and survey reefs and wrecks. A bounding box or an
image-level label is too coarse for these uses: a robot must know which pixels
are open water it can move through, and a marine survey must measure how much
of the scene is covered by reef or sea-grass. We therefore study semantic
segmentation of underwater images.

The input is an RGB image $x\in\mathbb{R}^{3\times H\times W}$. The output is a
mask $\hat y\in\{0,\ldots,7\}^{H\times W}$ that assigns one of the eight SUIM
classes to every pixel, so the prediction unit is one pixel, not one image. A
model $f_\theta$ produces a score map $f_\theta(x)\in\mathbb{R}^{8\times H\times W}$,
and each pixel takes the class with the highest score,
\begin{equation}
\hat y_{uv}=\arg\max_{k\in\{0,\ldots,7\}} f_\theta(x)_{k,uv}.
\end{equation}
The classes are mutually exclusive and there is no void label, so every pixel of
every mask is a training target and is counted by the metrics.

Our research question has two parts. First, on SUIM, how much better is a
DeepLabV3 model pretrained on COCO images with the Pascal VOC categories than a
U-Net trained from scratch, measured by mIoU and per-class IoU on the official
test set? Second, with the DeepLabV3 model and all
other settings fixed, does adding a soft Dice term to the cross-entropy loss
raise IoU on the rare classes (RO, PF and WR) without lowering mIoU? The first
comparison changes both the architecture and the initialization, so it cannot
attribute the gap to pretraining alone; only the second comparison isolates a
single factor.

Four properties of underwater images make the task harder than segmentation of
ordinary photographs.
\begin{itemize}
\item \textit{Colour cast.} Water absorbs red light much faster than blue and
green light, so most images are tinted blue or green and the same object
changes colour with depth and distance. Colour is therefore a weak cue for
classes such as reef, rock and sea-floor.
\item \textit{Turbidity.} Suspended particles scatter light, which lowers
contrast and blurs edges. Object boundaries become uncertain both for the model
and for the human annotators.
\item \textit{Small objects.} Divers, robots, fish and plants usually cover a
small part of the frame (Section~\ref{sec:eda}), so a large share of their
pixels lie on a boundary and they lose detail when the image is resized.
\item \textit{Class imbalance.} Reefs and the waterbody background cover most
labelled pixels, while robots cover only \suimShareTrainValRO\% of them. A model
can reach high pixel accuracy while ignoring the rare classes, which is why we
report mIoU, which weights every class equally.
\end{itemize}

\subsection{Dataset and study scope}
\label{sec:dataset}
We use SUIM (Segmentation of Underwater IMagery), published by the Interactive
Robotics and Vision Laboratory of the University of Minnesota together with the
IROS 2020 paper of Islam et al.\ \cite{suim}, and downloaded from the dataset page
\cite{suim-data}. Neither the paper nor the information file in the archive gives
a version number, so we identify our copy by its download date, 02 Oct 2026.
According to the authors, the images were taken during ocean explorations and
human-robot cooperative experiments at several sites with different water types,
and a few come from their earlier underwater datasets EUVP, USR-248 and UFO-120.
Every image was annotated by hand with one dense pixel-level mask, and the
annotation was done by seven human participants \cite{suim}.

Each mask stores its class as a 3-bit RGB colour in which each channel is either
0 or 255. Read as a binary number, the colour gives the class index: background
waterbody (BW, 000, black), human divers (HD, 001, blue), aquatic plants and
sea-grass (PF, 010, green), wrecks and ruins (WR, 011, sky blue), robots and
instruments (RO, 100, red), reefs and invertebrates (RI, 101, pink), fish and
vertebrates (FV, 110, yellow) and sea-floor and rocks (SR, 111, white). BW is the
background class; the other \suimForegroundClasses{} are foreground classes.
Images are JPEG files and masks are 24-bit BMP files, paired by filename. The
archive holds \suimOnDiskTotal{} pairs, \suimTrainValOnDisk{} in the training
pool and \suimTestOnDisk{} in the official test set, of which \suimUsableTotal{}
remain usable after the audit in Section~\ref{sec:eda}. Image sizes vary, with
\suimResolutionCount{} native resolutions. All pairs together occupy
\suimStorage{}~GiB, of which the images take \suimImageStorage{}~GiB and the masks
\suimMaskStorage{}~GiB, because the masks are stored as uncompressed bitmaps while
the images are compressed JPEG files.

The authors' code repository is released under the MIT licence, but that licence
covers the software, and the dataset page states no separate licence for the
images and masks. We therefore do not assume that the dataset may be
redistributed: we use it only for this course project, cite its source, and keep
it out of our repository, which provides a download script instead. We use only
the image and mask pairs. The archive also contains the authors' pretrained
checkpoints, their benchmark predictions on the test set, and extra per-class
binary masks and saliency masks for the test images; we use none of them, so all
models in this report are trained by us.

The dataset has three limitations that affect how far our results generalize.
Because SUIM reuses images from EUVP, USR-248 and UFO-120, its scenes share the
sites, cameras and water conditions of those datasets, and the test set may
contain scenes related to training scenes without being exact copies
(Section~\ref{sec:splits}). The data come from a limited set of locations, so
scores on SUIM may not carry over to other waters. Finally, each image has a
single annotation and the masks contain no void label for ambiguous regions such
as blurred edges, so some boundary errors reflect label uncertainty rather than
model mistakes.

\subsection{Dataset proposal and approval}
\label{sec:proposal}
The Dataset Proposal submitted at Milestone 1 fixes the track (semantic
segmentation), the dataset (SUIM), the split protocol, the metrics (mIoU, mDice and per-class
IoU), the two baselines, the pretrained model, the controlled loss experiment and
the compute estimate. Its final revision was committed on 07 Oct 2026, the
Milestone 1 deadline.

\textbf{Approval status:} Approved. \textbf{Conditions:} none.

SUIM meets every default segmentation requirement of Handbook Section 19.4
without an exception. It has \suimForegroundClasses{} foreground classes plus
the background class, against a minimum of three. It has
\suimUsableTotal{} usable image and mask pairs after the audit in
Section~\ref{sec:eda}, against a threshold of about 1,000. The per-class pixel
distribution required by the handbook is analysed in Section~\ref{sec:eda}.

\subsection{Exploratory data analysis}
\label{sec:eda}
Unless stated otherwise, the statistics in this section are computed on the
\suimTrainValUsable{} usable pairs of the official training pool. These pairs
remain after removing the masks rejected under mask quality below and the
copies of test images described in Section~\ref{sec:splits}. The official test set
is measured separately and only to describe it.

\paragraph{Pixel distribution per class.}
Table~\ref{tab:eda-classes} and Figure~\ref{fig:class-distribution} count each
class in two ways: the number of images that contain it and its share of all
labelled pixels. Reefs and invertebrates (RI) and the waterbody background (BW)
cover \suimShareTrainValRI\% and \suimShareTrainValBW\% of the pixels
respectively, while robots (RO) cover only \suimShareTrainValRO\% and appear in
\suimImagesTrainValRO{} images. The largest foreground class therefore has about
\suimImbalanceRatio{} times as many pixels as the smallest. Fish and vertebrates
(FV) show the opposite pattern: they appear in \suimImagesTrainValFV{} images
but cover only \suimShareTrainValFV\% of the pixels, because each fish is small.
This matters for evaluation and training. A model that ignores RO, PF and HD
loses little pixel accuracy, but mIoU weights all eight classes equally, so those
classes count as much as RI. The same imbalance motivates the cross-entropy
versus cross-entropy plus Dice experiment in Section~\ref{sec:experiments}.

The last column of Table~\ref{tab:eda-classes} shows that the official test set
has a different class balance from the training pool. BW rises from
\suimShareTrainValBW\% to \suimShareTestBW\% of the pixels, and RI falls from
\suimShareTrainValRI\% to \suimShareTestRI\%. Test scores may therefore differ
from validation scores because of this shift, not only because of overfitting.

% Generated by scripts/report/a2_data_tables.py from results/a2/eda/suim/summary.json. Do not edit.
\begin{table}[htbp]
\centering\small
\caption{Per-class statistics. Images, pixel share and median area are
computed on the usable train and validation pool; median area is the class
fraction of an image, over images that contain the class. The last column is
the pixel share on the official test set.}
\label{tab:eda-classes}
\begin{tabular}{llrrrr}
\toprule
Code & Class & Images & Pixels (\%) & Median area (\%) & Test pixels (\%)\\
\midrule
BW & Background (waterbody) & 1{,}185 & 30.99 & 33.6 & 41.81\\
HD & Human divers & 359 & 2.34 & 4.9 & 3.69\\
PF & Aquatic plants and sea-grass & 218 & 2.29 & 7.4 & 3.08\\
WR & Wrecks and ruins & 245 & 6.77 & 43.2 & 7.28\\
RO & Robots and instruments & 86 & 0.56 & 3.8 & 0.90\\
RI & Reefs and invertebrates & 964 & 35.64 & 56.1 & 19.47\\
FV & Fish and vertebrates & 958 & 7.26 & 7.2 & 6.77\\
SR & Sea-floor and rocks & 571 & 14.15 & 29.5 & 17.00\\
\bottomrule
\end{tabular}
\end{table}


\paragraph{Region size and image resolution.}
Figure~\ref{fig:region-area} shows how much of an image a class occupies when it
is present. RI and wrecks and ruins (WR) usually fill a large part of the frame,
with median areas of \suimMedianAreaRI\% and \suimMedianAreaWR\%. RO, HD, FV and
PF are small objects, with median areas between \suimMedianAreaRO\% and
\suimMedianAreaPF\%. Small regions are affected most by downsampling to the
training resolution, and a larger share of their pixels lies on a boundary. We
therefore expect boundary errors to dominate on these classes and report
boundary and interior IoU separately in the error analysis.

The usable pairs come in \suimResolutionCount{} native resolutions
(Figure~\ref{fig:resolutions}). About \suimDominantResolutionShare\% are
\suimDominantResolution, so every model receives the same fixed 4:3 input size
described in Section~\ref{sec:preprocessing}. Images with other aspect ratios,
mostly 16:9, are stretched by this resize; this is a known limitation.

\reportfigure{Images/suim_class_distribution.png}{Pixel share and image count
for each SUIM class, computed on the usable train and validation pairs.}{fig:class-distribution}

\reportfigure{Images/suim_region_area_distribution.png}{Distribution of region
area per class, as a fraction of the image, for images that contain the
class.}{fig:region-area}

\reportfigure{Images/suim_resolution_distribution.png}{Distribution of native
image resolutions in the train and validation pool.}{fig:resolutions}

\paragraph{Mask quality.}
The archive contains \suimOnDiskTotal{} image and mask pairs:
\suimTrainValOnDisk{} in the training pool and \suimTestOnDisk{} in the test set.
Images and masks are matched by filename stem. The number of unmatched files is
\suimUnpaired{}, and the number of files that fail to decode is
\suimDecodeErrors{}. The audit then checks
every mask pixel against the eight exact class colours, in which each RGB channel
is either 0 or 255.

In the training pool, \suimOffCodeMasks{} masks contain pixels with other
colours, while the test set has \suimTestOffCodeMasks{} such masks. We reject a
mask if even one of its pixels is off code. This strict rule does not depend on
the cut-off we choose: every other mask has no off-code pixel at all, while the
rejected masks have between \suimOffCodeShareMin\% and \suimOffCodeShareMax\% of
their pixels off code, so any tolerance below \suimOffCodeShareMin\% would reject
exactly the same masks. The
\suimSizeMismatches{} pairs whose mask and image have different sizes are all
among these \suimOffCodeMasks{} masks. Figure~\ref{fig:off-code-examples} shows
one rejected mask of each kind. Both look clean to the eye, yet
\suimOffCodeShareSameSize\% and \suimOffCodeShareTaller\% of their pixels are off
code, and these pixels form square blocks along object edges, the pattern typical
of JPEG compression. The second mask is also \suimTallerExtraRows{} rows taller
than its image, so its labels cannot be aligned with the image pixels.
Because the colour itself encodes the
class, an off-code pixel cannot be repaired reliably: thresholding each channel
maps it to the nearest code, which can create false rare-class pixels along
object edges. We therefore exclude the whole pair instead of repairing it, which
leaves \suimQualityPassed{} pairs in the training pool. Excluding these pairs
also removes some genuine rare-class examples, so RO results rest on fewer
images than the archive contains.

\reportfigure{Images/suim_off_code_examples.png}{Two masks rejected by the
audit. Off-code pixels are pixels whose colour is none of the eight class codes.
The examples are chosen by a fixed rule rather than by hand: among the rejected
masks with the same size as their image (top) and those taller than their image
(bottom), each row shows the mask with the median share of off-code
pixels.}{fig:off-code-examples}

\subsection{Splits and leakage prevention}
\label{sec:splits}
SUIM ships an official test set but no validation set. We keep all
\suimTestUsable{} official test pairs unchanged, so that test scores remain
comparable with published results. The test set was opened only for the audit
above. No test label or test score is used for hyperparameter tuning, early
stopping or checkpoint selection, and each selected model is evaluated on it
once, after its configuration is fixed.

\paragraph{Duplicates and leakage.}
SUIM contains repeated photographs, both inside the training pool and across the
training pool and the test set. We find them in two stages. First, each image is
reduced to a \suimHashBits-bit difference hash computed from a small greyscale
thumbnail; pairs whose hashes differ in at most \suimHashDistance{} bits become
candidates. This stage yields \suimHashCandidates{} candidate pairs. Second, a
candidate is confirmed only if the two images, resized to a common size, differ
by less than \suimPixelTolerance{} intensity levels per channel on average (on a
0 to 255 scale). This stage confirms \suimConfirmedPairs{} pairs. Confirmed
pairs are merged transitively into \suimDuplicateGroups{} duplicate groups.

Of these groups, \suimCrossSplitGroups{} contain both a training-pool image and a
test image. Figure~\ref{fig:duplicate-examples} shows the most and the least
similar of these pairs. The first is an exact copy, with an audit score of
\suimDuplicateScoreLowest{}. The second shows the same scene stored at a
different resolution and aspect ratio, so the files differ, yet its score of
\suimDuplicateScoreHighest{} is still far below the threshold of
\suimPixelTolerance{}. Every cross-split pair is therefore a copy of the same
photograph, not merely a similar scene. We keep the test images and drop their
\suimTestDuplicatesDropped{} copies from the training pool, which leaves
\suimTrainValUsable{} usable pairs for training and validation.

\reportfigure{Images/suim_duplicate_examples.png}{Training-pool copies of test
images, with the lowest (top) and highest (bottom) audit score among the
cross-split duplicate groups. The difference map is the absolute pixel
difference on a 0 to 255 scale, after resizing the training image to the test
image size. The audit score is the mean difference used to confirm a
duplicate.}{fig:duplicate-examples}

\paragraph{Split unit and stratification.}
The remaining \suimWithinGroups{} groups lie entirely in the training pool and
contain \suimWithinGroupImages{} images. The split unit is a duplicate group, not
an image, so every copy of a photograph lands on the same side of the split. The
\suimTrainValUsable{} pairs form \suimSplitGroups{} such units. To keep rare
classes on both sides, each image is labelled with the rarest class it contains,
and \texttt{StratifiedGroupKFold} with \suimValFolds{} folds, shuffling and seed
\suimSplitSeed{} assigns fold 0 to validation. Table~\ref{tab:data-splits} gives
the result: \suimTrainImages{} training and \suimValImages{} validation images,
a ratio of \suimTrainPercent{} to \suimValPercent{} percent. All eight classes
appear in both subsets, and no class differs by more than \suimMaxShareGap{}
percentage points in pixel share between them. RO, the rarest class, has
\suimImagesTrainRO{} training and \suimImagesValRO{} validation images.

% Generated by scripts/report/a2_data_tables.py from results/a2/eda/suim/summary.json. Do not edit.
\begin{table}[htbp]
\centering\small
\caption{SUIM splits. The official test set is held out; train and validation
come from the official training pool with seed 42.}
\label{tab:data-splits}
\begin{tabular}{lrrll}
\toprule
Split & Images & Share (\%) & Source & Use\\
\midrule
Train & 1{,}143 & 74.2 & Official train pool & Optimization\\
Validation & 287 & 18.6 & Official train pool & Checkpoint selection\\
Test & 110 & 7.1 & Official test set & Final evaluation only\\
\midrule
Total & 1{,}540 & 100.0 & & \\
\bottomrule
\end{tabular}
\end{table}


The membership of the training and validation sets is saved in two split files,
which are committed to the repository and have not changed since they were
created. The data loader reads these files instead of
re-splitting, and it stops with an error if any listed image is missing, so a
training run cannot silently use a different subset. Both training runs in
Section~\ref{sec:training-behavior} record \suimValImages{} validation images.

\paragraph{What has not been checked.}
Three sources of leakage or bias remain. First, the difference hash is not
invariant to mirroring, so a horizontally flipped copy of an image would not be
detected. Second, consecutive frames from the same video can differ by more than
the pixel tolerance while still showing the same scene; such frames may still
fall on both sides of a split. Third, SUIM reuses images from EUVP, USR-248 and
UFO-120, so related scenes may appear in both the training pool and the test set
without being pixel-level duplicates. In addition, the validation set contains
only \suimImagesValRO{} RO images, so validation IoU for RO is noisy and small
differences in it should not decide between models.

\subsection{Preprocessing and augmentation}
\label{sec:preprocessing}
The training and validation loaders read only the pairs listed in the split
files, so the pairs excluded from the training pool never reach training or
model selection. The test loader reads every official test pair, none of which
was excluded. Each mask is turned into a map
of class indices by thresholding every channel at \suimMaskThreshold{} and
reading the three bits as a binary number,
$c = 4\,[R \ge \suimMaskThreshold] + 2\,[G \ge \suimMaskThreshold] +
[B \ge \suimMaskThreshold]$, which gives 0 for BW (black) and 7 for SR (white).
Since every mask that passed the audit holds only the eight exact codes, the
threshold changes no pixel of our data; it only guards against a damaged file.

Every image is resized to $\suimInputHeight\times\suimInputWidth$ pixels
(height by width) with bilinear interpolation, and its mask with
nearest-neighbour interpolation, so that no new label appears on the boundary
between two classes. This size is half of the dominant resolution
\suimDominantResolution{} in each direction. Images with another aspect ratio
are stretched rather than cropped, so no labelled region is cut off. Validation
and test masks are resized in the same way, which means that all metrics in this
report are computed at $\suimInputHeight\times\suimInputWidth$, not at the
native resolution; thin structures and small regions can lose a few pixels in
this step. Pixel values are scaled to $[0, 1]$ and normalized with the ImageNet
channel mean \suimNormMean{} and standard deviation \suimNormStd{}. The
pretrained backbone expects these constants, and we apply them to U-Net as well
so that both models receive identical inputs.

The only augmentation is a random horizontal flip with probability one half,
applied to the image and its mask together and only during training. We add no
colour, crop or scale augmentation, so the colour cast and turbidity described
in Section~\ref{sec:task} reach the models as they are in the data. We also do
not reweight or resample classes; the only change aimed at class imbalance is
the Dice term in the controlled experiment.

A training batch therefore consists of an image tensor of shape
$B\times3\times\suimInputHeight\times\suimInputWidth$ (float32) and a mask
tensor of shape $B\times\suimInputHeight\times\suimInputWidth$ (int64) with
values from 0 to 7. The batch size $B$ is \suimUnetBatch{} for U-Net and
\suimDeeplabBatch{} for DeepLabV3; accumulating gradients over
\suimUnetAccumulation{} and \suimDeeplabAccumulation{} batches gives effective
batch sizes of \suimUnetEffectiveBatch{} and \suimDeeplabEffectiveBatch{}.
Figure~\ref{fig:example-batch} shows one batch of \suimExampleBatchSize{}
images drawn by the DeepLabV3 training loader. After normalization its values
lie between \suimExampleValueMin{} and \suimExampleValueMax{}, the flip was
applied to \suimExampleFlipped{} of the \suimExampleBatchSize{} images, and the
masks contain \suimExampleClasses{} of the eight classes. The overlay confirms
that the masks stay aligned with their images after resizing and flipping. The
batch also shows the limits of the labels: in the last image, a partly visible
animal at the right edge is labelled as sea floor.

\reportfigure{Images/suim_example_batch.png}{One training batch after
preprocessing, drawn by the DeepLabV3 training loader. Rows show the input after
undoing the normalization, the class mask in the SUIM colour code, and the mask
blended over the image, with background water left unblended. Each column title
gives the image name and whether the random flip was applied.}{fig:example-batch}
